\documentclass[12pt, a4paper]{article}
\usepackage{amsmath, amsthm, amssymb}
\usepackage{geometry}
\geometry{margin=2.5cm}

\newtheorem{theorem}{Theorem}
\newtheorem{remark}{Remark}

\title{\textbf{There Are Infinitely Many Prime Numbers}}
\author{0xHaru}
\date{\today}

\begin{document}
\maketitle

\begin{theorem}[Euclid's theorem]
The set of prime numbers is infinite.
\end{theorem}

\begin{proof}
Suppose for contradiction that there are only finitely many primes, say $k$ in
total. List them as $p_1, p_2, p_3, \ldots, p_k$, multiply them all together,
and consider the number that follows their product:
\[
    n = p_1 \times p_2 \times p_3 \times \cdots \times p_k + 1.
\]

\textbf{Case 1: $n$ is prime.} Then we have found a prime not in the list of all
$p_i$ --- a contradiction, since the list was supposed to be exhaustive.

\medskip

\textbf{Case 2: $n$ is composite.} By the Fundamental Theorem of Arithmetic,
$n$ must be a product of primes. However, none of the primes in the list can be
a divisor of $n$, because the remainder of every such division is $1$. It
follows that the prime factors of $n$ are primes outside the list --- again a
contradiction.

\medskip

Since both cases lead to a contradiction, the assumption that there are finitely
many primes is false. Therefore, the primes are infinite.
\end{proof}

\begin{remark}[Why the remainder is always $1$]
Think of it this way: the product $p_1 \times p_2 \times \cdots \times p_k$ is
a multiple of every $p_i$ in the list --- each one is literally baked into the
product as a factor. A multiple of $p_i$ is, by definition, something that
$p_i$ divides exactly, leaving no remainder. Now $n$ is obtained by adding $1$ to that product, stepping exactly one unit
past a number that $p_i$ divides cleanly. So no matter which $p_i$ you try
to divide by you always get a remainder of $1$.
\end{remark}

\end{document}
